\chapter{集合、函数与依赖关系}
\section{对象、元素与函数}
\begin{definition}[函数及其复合]
设 $A,B$ 为集合。函数 $f:A\to B$ 给每个 $a\in A$ 指定唯一的元素 $f(a)\in B$。若 $g:B\to C$，复合 $g\circ f:A\to C$ 由 $(g\circ f)(a)=g(f(a))$ 定义。恒等函数 $\id_A:A\to A$ 由 $\id_A(a)=a$ 定义。
\end{definition}
\begin{proposition}
函数复合满足结合律及单位律：
\[
h\circ(g\circ f)=(h\circ g)\circ f,
\qquad \id_B\circ f=f=f\circ\id_A.
\]
\end{proposition}
\begin{proof}
在任意 $a\in A$ 处，两边分别取值为 $h(g(f(a)))$，所以第一式成立。后两式在 $a$ 处的值均为 $f(a)$。这里使用了集合论中的函数外延原则，即两个定义域和陪域相同的函数相等，当且仅当它们逐点相等。
\end{proof}
\begin{definition}[纤维]
函数 $f:A\to B$ 在 $b\in B$ 上的纤维为
\[
f^{-1}(b)=\{a\in A\mid f(a)=b\}.
\]
函数 $f$ 为单射，是指每个纤维至多有一个元素；为满射，是指每个纤维至少有一个元素；为双射，是指每个纤维恰有一个元素。
\end{definition}
\begin{lemma}[双射的纤维判据]
函数 $f:A\to B$ 有双侧逆函数，当且仅当每个纤维恰有一个元素。
\end{lemma}
\begin{proof}
若 $g$ 为逆函数，则 $g(b)$ 属于 $f^{-1}(b)$。任意 $a\in f^{-1}(b)$ 满足 $a=g(f(a))=g(b)$。反之，取 $g(b)$ 为该纤维的唯一元素。唯一性保证这一定义不涉及任意选择。由定义 $f(g(b))=b$，而 $a$ 和 $g(f(a))$ 属于同一单元素纤维，故 $g(f(a))=a$。
\end{proof}
\begin{remark}
后文将把“单元素集合”替换为“可缩空间”。因此，“等价映射具有可缩的同伦纤维”直接推广了上述双射判据。理解这种替换时，必须保留纤维内的全部路径信息。
\end{remark}

\section{集合族与依赖和}
\begin{definition}[集合族]
以集合 $A$ 为指标的集合族记作 $(B_a)_{a\in A}$。其依赖和为
\[
\sum_{a\in A}B_a=\{(a,b)\mid a\in A,\ b\in B_a\}.
\]
它配备投影 $\pi(a,b)=a$。符号“依赖”表示第二个分量的取值范围可以随着第一个分量改变。
\end{definition}
\begin{example}
令 $A=\N$，$B_n=\{0,1,\ldots,n\}$，则
\[
\sum_{n\in\N} B_n=\{(n,k)\in\N^2\mid k\le n\}.
\]
若 $B_a=B$ 与 $a$ 无关，则依赖和等于笛卡尔积 $A\times B$。若每个 $B_a$ 为单元素集合，则其依赖和自然同构于 $A$。
\end{example}
\begin{proposition}[映射与纤维族]
每个函数 $p:E\to A$ 确定一个纤维族 $(E_a)_{a\in A}$，其中 $E_a=p^{-1}(a)$；存在与投影相容的自然双射
\[
E\cong\sum_{a\in A} E_a.
\]
\end{proposition}
\begin{proof}
定义 $\Phi(e)=(p(e),e)$。因为 $e\in E_{p(e)}$，这是合法的有序对。逆映射为 $\Psi(a,e)=e$。有 $\Psi\Phi(e)=e$；若 $e\in E_a$，则 $p(e)=a$，所以 $\Phi\Psi(a,e)=(a,e)$。此外，两侧投影在 $e$ 上的值均为 $p(e)$。
\end{proof}
\begin{definition}[族之间的映射]
两个 $A$-指标族 $(B_a)$、$(C_a)$ 之间的逐纤维映射为一族函数 $u_a:B_a\to C_a$。它对应于一个总空间映射
\[
u:\sum_a B_a\longrightarrow\sum_a C_a,
\qquad u(a,b)=(a,u_a(b)),
\]
且与到 $A$ 的投影交换。
\end{definition}
\begin{remark}
“在同一底空间之上”会成为一种重要约束。它要求保持参数 $a$，并允许在该参数对应的纤维内部变化。若没有这一约束，总空间映射可以把不同参数的纤维混在一起。
\end{remark}

\section{依赖积与截面}
\begin{definition}[依赖积]
集合族 $(B_a)_{a\in A}$ 的依赖积为
\[
\prod_{a\in A}B_a
=\{s\mid s\text{ 是函数，且对每个 }a,\ s(a)\in B_a\}.
\]
这样的 $s$ 称为依赖函数。若所有 $B_a$ 都等于 $B$，则依赖积为函数集合 $B^A$。
\end{definition}
\begin{definition}[截面]
映射 $p:E\to A$ 的截面为函数 $s:A\to E$，满足 $p\circ s=\id_A$。
\end{definition}
\begin{proposition}
依赖积 $\prod_a B_a$ 自然等同于投影 $\sum_a B_a\to A$ 的截面集合。
\end{proposition}
\begin{proof}
依赖函数 $s$ 给出截面 $a\mapsto(a,s(a))$。反之，任一截面的第一坐标必为 $a$，故有唯一表达式 $a\mapsto(a,s(a))$，其中 $s(a)\in B_a$。两种构造互逆。
\end{proof}
\begin{example}[空指标]
空依赖和为 $\sum_{a\in\varnothing}B_a=\varnothing$。空依赖积则是单元素集合，因为从空集合出发恰有一个函数。这个差别将在“存在量词”和“全称量词”的解释中反复出现。
\end{example}
\begin{proposition}[依赖柯里化]
设 $C_{a,b}$ 为依赖于 $a\in A$、$b\in B_a$ 的集合族。存在双射
\[
\prod_{(a,b)\in\sum_aB_a}C_{a,b}
\cong \prod_{a\in A}\prod_{b\in B_a}C_{a,b}.
\]
\end{proposition}
\begin{proof}
左侧函数 $h$ 对应于右侧函数 $a\mapsto(b\mapsto h(a,b))$。逆方向把 $k$ 送到 $(a,b)\mapsto k(a)(b)$。逐点代入验证两者互逆。参数之间的依赖顺序没有改变，故所有表达式都具有确定的定义域。
\end{proof}
\begin{proposition}[依赖选择的有序对公式]
存在自然双射
\[
\prod_{a\in A}\sum_{b\in B_a}C_{a,b}
\cong
\sum_{f\in\prod_a B_a}\prod_{a\in A} C_{a,f(a)}.
\]
\end{proposition}
\begin{proof}
若 $u(a)=(b_a,c_a)$，令 $f(a)=b_a$，$v(a)=c_a$，得到 $(f,v)$。反之，由 $(f,v)$ 得到 $a\mapsto(f(a),v(a))$。两次拆分和重新组合后，各个坐标保持不变。这一双射本身不需要选择公理，因为左边已给出对每个 $a$ 的具体有序对。
\end{proof}
\begin{remark}
必须区别“每个纤维非空”与“给定一个截面”。在经典集合论中，一般非空集合族存在截面是选择公理的内容；上述命题则只是已有数据的重组。类型论中的 $\sum$ 会记录见证，因此也必须保持这一区别。
\end{remark}

\section{代入、拉回与交换图}
\begin{definition}[拉回]
给定 $f:X\to A$、$p:E\to A$，定义
\[
X\times_AE=\{(x,e)\in X\times E\mid f(x)=p(e)\}.
\]
其投影组成交换方块
\[
\begin{tikzcd}
X\times_AE\ar[r,"\pi_E"]\ar[d,"\pi_X"']&E\ar[d,"p"]\\
X\ar[r,"f"']&A.
\end{tikzcd}
\]
\end{definition}
\begin{proposition}[拉回的泛性质]
若 $u:T\to X$、$v:T\to E$ 满足 $fu=pv$，则存在唯一 $w:T\to X\times_AE$，使得 $\pi_Xw=u$、$\pi_Ew=v$。
\end{proposition}
\begin{proof}
只能定义 $w(t)=(u(t),v(t))$。条件 $fu=pv$ 保证这个点属于拉回。两个投影条件成立；任一满足投影条件的映射都具有相同的两个坐标，故唯一。
\end{proof}
\begin{proposition}[拉回是参数代入]
若 $E=\sum_{a\in A}B_a$，则作为 $X$ 上的族，$X\times_AE$ 自然同构于
\[
\sum_{x\in X}B_{f(x)}.
\]
\end{proposition}
\begin{proof}
拉回中的元素形如 $(x,(a,b))$，且 $f(x)=a$。删去由 $x$ 确定的 $a$，得到 $(x,b)$，其中 $b\in B_{f(x)}$。逆方向插入 $a=f(x)$。两种映射都保持到 $X$ 的投影。
\end{proof}
\begin{example}
把 $B_n=\R^n$ 沿 $f:\N\to\N$、$f(k)=2k$ 拉回，得到 $(\R^{2k})_{k\in\N}$。这里没有增加新的纤维运算，只改变了参数如何进入原有的族。
\end{example}

\section{命题与证明数据}
\begin{definition}[命题的集合解释]
一个命题可以用空集合或单元素集合表示，分别对应假和真。若 $P(a)$ 是关于 $a\in A$ 的命题，令 $B_a$ 为其真值集合，则 $\prod_aB_a$ 表示对所有 $a$ 成立，而 $\sum_aB_a$ 记录一个使 $P(a)$ 成立的 $a$。
\end{definition}
\begin{remark}
依赖和可能有多个元素。因此它记录的内容多于单纯的“存在”。例如 $\sum_{n\in\N}\{\text{$n$ 为偶数的证明}\}$ 还记录具体偶数。后文的命题截断将保留存在性并遗忘具体见证。
\end{remark}
\begin{proposition}[单元素纤维的截面]
如果每个 $B_a$ 恰有一个元素，则 $\prod_aB_a$ 恰有一个元素。
\end{proposition}
\begin{proof}
给 $a$ 指定 $B_a$ 的唯一元素，得到一个依赖函数。任意两个这样的函数在每个 $a$ 处取值相同，故由函数外延性相等。
\end{proof}
\begin{remark}
上述命题的同伦版本是“可缩空间族的依赖积可缩”。它会保证很多构造在全部参数中同时成立，而不仅是在逐个参数上存在。由具体集合过渡到空间时，这种整体参数依赖是最需要注意的部分。
\end{remark}
